\begin{textsample}{NEG Dim 3 – 2023\_c5 – Score 3.00 – 2023\_c5/enade\_2023\_tec\_seg\_trab\_4f8342634e.txt}  \label{ex:f3_neg_229}
CAPÍTULO 4 DISTRIBUIÇÃO DE CONCEITOS A finalidade deste capítulo é apresentar um panorama nacional sobre a distribuição do Conceito Enade 2023.
Para traçar esse panorama, recorre -se a tabelas e gráficos, que são acompanhados de comentários e análises que contribuem para a contextualização dos dados neles expostos.
Para visualizar o panorama do Enade 2023, em nível nacional, na primeira seção deste capítulo, os dados sobre o número de cursos participantes são organizados por grande região.
Na segunda, o mapeamento do conceito do exame é feito por categoria administrativa das IES e modalidade de oferta dos cursos, de maneira cruzada às grandes regiões.
Na última seção, a lógica de apresentação dos dados é feita por organização acadêmica e, também, por grande região de funcionamento dos cursos, no país como um todo. 4.1 PANORAMA NACIONAL DA DISTRIBUIÇÃO DOS CONCEITOS A lógica adotada na apresentação dos dados, nesta seção e nas duas seguintes, pressupõe o registro de duas informações de natureza reiterativa.
O primeiro diz respeito aos valores possíveis do Enade, que variam em níveis de 1 a 5, sendo que, à medida que esses valores aumentam, o curso é melhor avaliado, e ao fato de que aqueles cursos que tiveram somente um ou nenhum concluinte participante não obtêm o Conceito Enade, ficando sem conceito ( SC ), como ressaltado na apresentação deste relatório.
O segundo registro é que, caso haja diferença entre os cursos tabulados neste capítulo e no capítulo 1, ela corresponderá aos cursos que não foram avaliados, ou seja, àqueles, em princípio, sem estudantes concluintes que tivessem sido inscritos para a prova.
Na Tabela 4.1, são apresentadas a quantidade e a distribuição de cursos de Tecnologia em Segurança no Trabalho participantes do Enade 2023, por faixa de conceito e grande região.
A partir dos dados da Tabela 4.1, nota -se que, dos 24 cursos participantes, 6 ( 25,0 \% ) classificaram -se com conceito 2, o valor modal.
Esse foi também o conceito modal nas regiões : Sudeste ( 42,9 \% ), Sul ( 25,0 \% ).
O conceito 3 foi o segundo mais frequente em nível nacional ( 20,8 \%, correspondendo a 5 cursos ), tendo sido o segundo conceito mais frequente na maioria das regiões, exceto no Nordeste e Norte.
O conceito 4 foi o terceiro ( 12,5 \%, correspondendo a 3 cursos ).
Houve, ainda, 3 ( 12,5 \% ) cursos que receberam o conceito 1 e dois ( 8,3 \% ) cursos que receberam o conceito 5.
Dos 24 cursos de Tecnologia em Segurança no Trabalho, 5 ( 20,8 \% ) ficaram sem conceito ( SC ).
A região Norte participou com 1 cursos, o que corresponde a 4,2 \% do total nacional.
Desses, um curso recebeu o conceito 5, o que equivale a 100,0 \% do total regional, o conceito modal na região.
A região Nordeste participou com 8 cursos, correspondentes a 33,3 \% do total nacional.
Nessa região, um curso ( 12,5 \% em termos regionais ) obteve conceito 5,.
Os conceitos 4 e 2 foram atribuídos, respectivamente, a um curso cada ( 12,5 \% e 12,5 \%, respectivamente ).
Ainda nessa região, um curso ( 12,5 \% ) foi avaliado com o conceito 1 e nenhum curso ( 0,0 \% ) foram avaliados com o conceito 3.
Outros quatro cursos ( 50,0 \% ) ficaram sem conceito ( SC ), o conceito modal para a região.
Para os 7 cursos participantes da região Sudeste ( 29,2 \% do total nacional ), o conceito 3 junto ao conceito 2, foram obtidos por três cursos ( 42,9 \% ).
O conceito 1, a um curso ( 14,3 \% ).
Receberam os conceitos 5 e 4, respectivamente, nenhum curso ( 0,0 \% ) e nenhum curso ( 0,0 \% ).
Nessa região, nenhum curso ( 0,0 \% ) ficou sem conceito ( SC ).
Os 8 cursos da região Sul corresponderam a 33,3 \% do total nacional.
Os conceitos 2, 3 e 4 foram atribuídos, cada, a 25,0 \% dos cursos da região, ou seja, atribuídos a 2 dos 8 cursos participantes na região Sul.
O conceito 1, a um curso ( 12,5 \% ).
Na região Sul, nenhum curso ( 0,0 \% ) obteve conceito 5 e um curso ( 12,5 \% ) ficou sem conceito ( SC ).
A região Centro-Oeste não participou com nenhum curso, por isso a ausência de dados.
Em um gráfico de distribuição cumulativa, em cada valor de abscissa ( eixo x ), o valor apresentado no eixo das ordenadas ( eixo y ) é o valor acumulado da variável ( Conceito Enade ).
Em linhas gerais, em um gráfico desta mesma tipologia, poligonais mais à esquerda representam, no caso dos dados apresentados nesta seção, as grandes regiões com a pior distribuição, e poligonais mais à direita, as grandes regiões com os melhores conceitos.
A partir destas diretrizes, no Gráfico 4.1, é apresentada a distribuição cumulativa do Conceito Enade, segundo a grande região.
No Gráfico 4.1, observando a linha vertical pontilhada, que ajuda a visualizar o percentual de conceitos acumulados até o valor 2, ou seja, conceitos 1 e 2, pode -se notar, por exemplo, que a região Sudeste ( linha vermelha ) apresenta 57,1 \% dos cursos com conceito 2 ou abaixo deste valor.
A região Norte ( linha azul ), por outro lado, apresenta 0,0 \% dos cursos na mesma situação : com conceito 2 ou abaixo.
Destaca -se que, a região Sudeste ( linha vermelha ), os conceitos 4 e 5 não foram alcançados por nenhum curso, na qual todos os cursos têm conceito menor ou igual a 3, e a respectiva poligonal cumulativas alcançam 100 \% no conceito 4.
A região Sudeste ( linha vermelha ) apresenta os cursos com conceitos mais baixos, ao passo que a região Norte ( linha azul ) apresenta a distribuição com valores maiores.
Os cursos em IES situadas na região Nordeste ( linha verde ) e Sul ( linha roxa ) apresentam uma situação intermediária. 4.2 CONCEITOS POR CATEGORIA ADMINISTRATIVA, MODALIDADE DE OFERTA DOS CURSOS E GRANDE REGIÃO Na Tabela 4.2, é apresentada a distribuição dos cursos de Tecnologia em Segurança no Trabalho participantes do Enade 2023, por categoria administrativa e modalidade de oferta, de acordo com os conceitos por eles alcançados, segundo as grandes regiões brasileiras.
De acordo com as informações da Tabela 4.2, dos 24 cursos participantes, 2 ( 8,3 \% ) eram ministrados em instituições públicas e 22 ( 91,7 \% ), em instituições privadas.
Quanto à modalidade de oferta, 5 ( 20,8 \% ) cursos eram presenciais e 19 ( 79,2 \% ), a distância.
De acordo com as informações da mesma tabela, em termos nacionais, entre instituições públicas e privadas, observa -se que, dos 2 cursos avaliados com conceito 5, 1 foram oferecidos em IES públicas, e 1, em IES privadas.
Dos 22 cursos participantes de IES privadas, o conceito 2 foi o valor modal, atribuído a 6 ( 27,3 \% ) cursos.
Em relação aos demais cursos de instituições privadas, cinco cursos ( 22,7 \% ) foram avaliados com o conceito 3, três cursos ( 13,6 \% ) foram avaliados com o conceito 1, dois cursos ( 9,1 \% ) foram avaliados com o conceito 4 e, um curso ( 4,5 \% ), com o conceito 5.
Nessa categoria, 5 cursos ( 22,7 \% ) ficaram sem conceito ( SC ).
Na rede pública, os conceitos 4 e 5 foram atribuídos a 1 curso, correspondente a 50,0 \% dos 2 cursos da categoria.
Dos 5 cursos participantes da modalidade presencial, o conceito 5 foi o valor modal, atribuído a 2 ( 40,0 \% ) cursos.
Dos demais cursos dessa modalidade de oferta, um curso ( 20,0 \% ) recebeu o conceito 4 ; nenhum curso ( 0,0 \% ) recebeu o conceito 3 ; nenhum curso ( 0,0 \% ), o conceito 2 ; e nenhum curso ( 0,0 \% ), o conceito 1.
Nessa modalidade de oferta, dois cursos ( 40,0 \% ) ficaram sem conceito ( SC ).
Dos 19 cursos participantes da modalidade a distância, o conceito 2 foi o valor modal, atribuído a 6 ( 31,6 \% ) cursos.
Dos demais cursos ofertados nessa modalidade, cinco cursos ( 26,3 \% ) receberam conceito 3, três cursos ( 15,8 \% ) receberam o conceito 1 ; dois cursos ( 10,5 \% ) receberam o conceito 4 ; e nenhum curso ( 0,0 \% ) recebeu conceito 5.
Nessa modalidade de oferta, três cursos ( 15,8 \% ) ficaram sem conceito ( SC ).
Na análise por região, observa -se, que, na região Norte, as instituições públicas participaram com nenhum curso ( 0,0 \% total regional ).
As instituições privadas participaram com 1 cursos na região Norte ( 100,0 \% do total regional ), o qual o conceito 5 foi o conceito modal nas instituições privadas na região, atribuído a um curso, o que corresponde a 100,0 \% dessa categoria na região.
Dos cursos oferecidos na região Norte, um era presencial e nenhum ficou sem conceito ( SC ).
Do único curso oferecido, nessa modalidade, obteve o conceito 5 ( modal ) ; Nenhum curso ficou sem conceito ( SC ) ou recebeu os demais conceitos na região Norte e na modalidade a distância.
Na região Nordeste, a rede privada concentra 6 dos 8 cursos participantes, o equivalente a 75,0 \% do total da região.
Entre os cursos oferecidos por IES privadas no Nordeste, um curso foi avaliado com conceito 2, o conceito modal.
Um curso obteve o conceito 1 ; nenhum curso, o conceito 5 ; nenhum curso, o conceito 4 ; e nenhum curso obteve o conceito 3.
Nessa categoria administrativa e grande região, quatro cursos ficaram sem conceito ( SC ).
As instituições públicas da região Nordeste participaram com dois cursos ( 25,0 \% do total da região ).
Desses, um ( 50,0 \% ) obteve conceito 5, o valor modal, seguidos por um curso, que recebeu o conceito 4.
O conceito 3 foi atribuído a nenhum curso ; o conceito 2, a nenhum curso ; e o conceito 1, a nenhum curso.
Nenhum curso ficou sem conceito ( SC ).
Dos cursos oferecidos na região Nordeste ( oito ), ( quatro ) eram presenciais : dois ficaram sem conceito ( SC ), nenhum obteve o conceito 1 ; nenhum, o conceito 2 ; nenhum, o conceito 3 ; um, o conceito 4 ; e um, o conceito 5.
Outros quatro cursos eram na modalidade a distância, tendo um curso recebido o conceito 1 ; um, o conceito 2 ; nenhum, o conceito 3 ; e nenhum, o conceito 5.
Dois cursos ficou sem conceito ( SC ).
Na região Sudeste, a proporção de cursos da rede privada, 100,0 \%, correspondeu a 7 dos 7 cursos participantes.
Nessa combinação de categoria administrativa e grande região, o 2 foi atribuído a três cursos ; e nenhum curso ficou sem conceito ( SC ).
Os demais foram avaliados com os seguintes conceitos : 2 ( três cursos ), 1 ( um curso ), 5 ( nenhum curso ) e 4 ( nenhum curso ).
Dos nenhum cursos oferecidos por instituições públicas na região Sudeste, nenhum obteve o conceito 5 ; nenhum, o conceito 4 ; nenhum, o conceito 3 ; nenhum obteve conceito 2 e nenhum obteve o conceito 1.
Nenhum curso ficou sem conceito ( SC ).
Ainda na região Sudeste, os cursos ( 0,0 \% ) foram oferecidos na modalidade a distância e nenhum na modalidade presencial : As instituições privadas concentraram 8 dos 8 cursos participantes da região Sul, 100,0 \% do total regional.
Nessa combinação de categoria administrativa e grande região, dois receberam o conceito 3 ; dois receberam o conceito 2 ; dois receberam o conceito 4 ; um, o conceito 1 ; e nenhum obteve o conceito 5.
Um curso ficou sem conceito ( SC ).
As instituições públicas na região Sul participaram com nenhum curso.
Na região Sul, na modalidade a distância, foram oferecidos oito cursos, sendo um curso com conceito 1, dois cursos com conceito 2, dois cursos com conceito 3 ( modal ), dois cursos com conceito 4 e nenhum curso com conceito 5.
Um curso ficou sem conceito ( SC ) nessa região.
Na região Centro-Oeste, não houve dados para categorias administrativas, bem ccomo modalidade de oferta, segundo a região e o Conceito Enade.
A título de síntese, no Gráfico 4.2, apresenta -se a distribuição do Conceito Enade, segundo a categoria administrativa da IES.
Nesse gráfico, poligonais mais à esquerda representam Conceito Enade mais baixos, e poligonais mais à direita, Conceito Enade mais altos.
É importante observar que, para cada poligonal presente nesse gráfico, a porcentagem constante no eixo y refere -se ao universo de cursos oferecidos por instituições públicas ( poligonal azul ) ou privadas ( poligonal verde ), respectivamente.
Os cursos em IES públicas ( poligonal azul ) apresentam uma poligonal mais à direita do que os em IES privadas ( poligonal verde ) e, consequentemente, uma distribuição de Conceito Enade com valores maiores. 4.3 CONCEITOS POR ORGANIZAÇÃO ACADÊMICA E GRANDE REGIÃO Na Tabela 4.3, encontra -se a distribuição dos conceitos atribuídos aos cursos participantes do Enade 2023, na área de Tecnologia em Segurança no Trabalho, por organização acadêmica, segundo as grandes regiões brasileiras.
Como mostram os dados da Tabela 4.3, dos 24 cursos de Tecnologia em Segurança no Trabalho participantes, 11 foram oferecidos em universidades, 9 em centros universitários e 3 em faculdades ; um curso foi oferecido por CEFET/IF.
Essa distribuição corresponde a, respectivamente, 45,8 \%, 37,5 \%, 12,5 \% e 4,2 \% dos cursos.
De acordo com os dados apresentados na mesma tabela, de todos os 5 cursos avaliados com o conceito 3, 3 eram vinculados a universidades.
Os cursos vinculados a esse tipo de organização acadêmica tiveram o conceito 2 como conceito modal, atribuído a 4 cursos ( 36,4 \% ).
Os demais cursos ofertados por universidades receberam os conceitos 1 ( dois cursos ), 5 ( um curso ) e 4 ( nenhum curso ).
Um curso ficou sem conceito ( SC ).
Dos cursos vinculados a centros universitários, o conceito modal foi o 3, atribuído a 2 cursos ( 22,2 \% ).
Três cursos vinculados a esse tipo de organização acadêmica ficaram sem conceito ( SC ), e os outros cursos receberam os conceitos 5 ( um curso ), 4 ( um curso ), 2 ( um curso ) e 1 ( um curso ).
Dos 3 cursos mantidos por faculdades, 1 ( 33,3 \% ) recebeu conceito 4, o conceito modal.
O segundo conceito mais frequentemente atribuído aos cursos oferecidos por essa categoria de IES foi o conceito 4, recebido por um curso.
Na sequência, nenhum curso obteve conceito 5 ; nenhum curso, o conceito 3 ; e nenhum curso, o conceito 1.
Dos cursos oferecidos em faculdades, um curso ficaram sem conceito ( SC ).
Dos cursos oferecidos por CEFET/IF, o conceito modal foi o 4, atribuído a um curso ( 100,0 \% ).
Considerando -se, separadamente, as regiões brasileiras, verifica -se que, na região Norte, os centros universitários foram representados por um curso, o qual foi atribuído ao conceito 5, conceito modal. por isso a ausência de dados atribuídos aos Conceito Enade.
As universidade, as faculdades e os CEFET/IF participaram com nenhum curso na região Norte.
Na região Nordeste, as universidades participaram com 4 dos 8 cursos da área de Tecnologia em Segurança no Trabalho oferecidos.
O conceito 5 foi atribuído a um curso.
Outro um curso obteve o conceito 2 ; um curso, o conceito 1 ; nenhum curso, o conceito 4 ; e nenhum curso, o conceito 3.
Um curso ficou sem conceito ( SC ).
Os centros universitários contaram com dois cursos participantes na região Nordeste, e esses cursos oferecidos por centros universitários no Nordeste ficaram sem conceito ( SC ).
As faculdades foram representadas por um curso na região Nordeste, o qual ficou ficou sem conceito ( SC ).
Na região Nordeste, os CEFET/IF participaram com um curso, o qual foi atribuído o conceito 4.
Na região Sudeste, as universidades concentraram 4 dos 7 cursos de Tecnologia em Segurança no Trabalho da região.
Dos cursos oferecidos em IES com este tipo de organização acadêmica, nessa região, o conceito modal foi 3, atribuído a dois cursos.
Os demais cursos receberam o conceito 2 ( dois cursos ).
Os centros universitários participaram com dois cursos na região Sudeste, dos quais um obteve o conceito modal 3.
O conceito 1 foi atribuído a um curso ; o conceito 5 a nenhum curso ; o conceito 4 a nenhum curso ; e o conceito 2 a nenhum curso.
Nenhum curso ficou sem conceito ( SC ) nessa combinação de organização acadêmica e região.
Já as faculdades foram representadas por um curso na região Sudeste, o qual nenhum ficou sem conceito ( SC ).
O conceito modal 2 foi atribuído a este curso.
Os CEFET/IF participaram com nenhum curso na região Sudeste.
Dos oito cursos da região Sul, três eram mantidos por universidades, para os quais o conceito 3 foi atribuído a um curso.
Os demais cursos receberam os conceitos 2 ( um curso ) e 1 ( um curso ).
Dos quatro cursos participantes de centros universitários da região Sul, um ficou sem conceito ( SC ), e um curso recebeu o conceito 4.
Os conceitos restantes ficaram assim distribuídos : 3 ( um curso ), 2 ( um curso ), 5 ( nenhum curso ) e 1 ( nenhum curso ).
Quanto aos um curso vinculados a faculdades na região Sul, um curso recebeu o conceito 4.
Na região Centro-Oeste, não houve cursos participantes vinculados às organizações acadêmicas.
De maneira semelhante à forma como se procedeu em relação à seção anterior, esta seção encerra -se com o Gráfico 4.3, que apresenta a distribuição do Conceito Enade dos cursos de Tecnologia em Segurança no Trabalho segundo a organização acadêmica da IES.
No Gráfico 4.3, apresenta -se a distribuição do Conceito Enade segundo a organização acadêmica da IES.
A leitura desse último gráfico é semelhante à dos anteriores, seguindo -se a lógica de que poligonais deslocadas para a esquerda correspondem aos conceitos mais baixos ( 1, 2 e 3 ) e poligonais deslocadas para a direita correspondem a Conceitos Enade mais altos.
Nas universidades ( linha rosa ), todos os cursos obtiveram conceito 5 ou abaixo disso, já que a poligonal cumulativa alcança 100 \% para o conceito 5.
Além disso, os cursos de CEFET/IF apresentam uma poligonal mais à direita do que as demais, denotando uma melhor distribuição de Conceitos Enade.
No outro extremo, encontram -se os cursos mantidos por universidades ( linha azul ), apresentando as respectivas poligonais mais à esquerda e, consequentemente, uma distribuição de Conceitos Enade com os menores valores.
Os cursos mantidos por faculdades ( linha vermelho ) e centros universitários ( linha verde ) apresentam poligonais mais centrais e, consequentemente, uma distribuição de Conceitos Enade com os valores intermediários.

% matched lemmas: 
\end{textsample}
